% Part: sets-functions-relations
% Chapter: relations
% Section: special-properties

\documentclass[../../../include/open-logic-section]{subfiles}

\begin{document}

\olfileid[fr]{sfr}{rel}{prp}
\olsection{Propriétés particulières des relations}

\begin{intro}
Certains types de relations sont si fréquents qu’on leur a donné des noms
particuliers. Ainsi, $\le$ et~$\subseteq$ mettent les objets de leurs domaines
respectifs (par exemple $\Nat$ pour~$\le$ et $\Pow{A}$ pour~$\subseteq$)
en relation de façons semblables. Pour préciser ce qui rapproche ces
relations et ce qui les distingue, nous les classons selon certaines
propriétés qu’elles peuvent posséder. Certaines combinaisons de ces
propriétés définissent des types de relations particulièrement importants :
les relations d’ordre et les relations d’équivalence.
\end{intro}

\begin{defn}[Réflexivité]
Une relation $R \subseteq A^2$ est \emph{réflexive} si et seulement si,
pour tout $x \in A$, on a $Rxx$.
\end{defn}

\begin{defn}[Transitivité]
Une relation $R \subseteq A^2$ est \emph{transitive} si et seulement si,
chaque fois que $Rxy$ et $Ryz$, on a aussi $Rxz$.
\end{defn}

\begin{defn}[Symétrie]
Une relation~$R \subseteq A^2$ est \emph{symétrique} si et seulement si,
chaque fois que $Rxy$, on a aussi~$Ryx$.
\end{defn}

\begin{defn}[Antisymétrie]
Une relation~$R \subseteq A^2$ est \emph{antisymétrique} si et seulement si,
chaque fois que l’on a à la fois $Rxy$ et $Ryx$, on a $x=y$
(autrement dit : si $x\neq y$, alors $\lnot Rxy$ ou $\lnot Ryx$).
\end{defn}

\begin{explain}
Pour une relation symétrique, $Rxy$ et $Ryx$ sont toujours vraies toutes
les deux ou fausses toutes les deux. Pour une relation antisymétrique,
$Rxy$ et $Ryx$ ne peuvent être vraies toutes les deux que si $x = y$.
Cela n’\emph{exige} pas que $Rxy$ et $Ryx$ soient vraies lorsque $x = y$ ;
cette possibilité n’est simplement pas exclue. Une relation antisymétrique
peut donc être réflexive, mais toute relation antisymétrique n’est pas
nécessairement réflexive. Remarquons aussi qu’être antisymétrique et ne pas
être symétrique sont deux conditions différentes. Une relation peut même
être à la fois symétrique et antisymétrique : c’est le cas, par exemple,
de la relation d’identité.
\end{explain}

\begin{defn}[Connexité]
Une relation $R \subseteq A^2$ est \emph{connexe} si, pour tous $x,y\in A$,
lorsque $x \neq y$, on a $Rxy$ ou~$Ryx$.
\end{defn}

\begin{prob}
Donnez des exemples de relations qui sont (a) réflexives et symétriques
mais non transitives, (b) réflexives et antisymétriques,
(c) antisymétriques et transitives mais non réflexives, et
(d) réflexives, symétriques et transitives. N’utilisez pas de relations
portant sur des nombres ou des ensembles.
\end{prob}

\begin{defn}[Irréflexivité]
Une relation $R \subseteq A^2$ est dite \emph{irréflexive} si,
pour tout $x \in A$, on n’a pas $Rxx$.
\end{defn}

\begin{defn}[Asymétrie]
Une relation $R \subseteq A^2$ est dite \emph{asymétrique} s’il n’existe
aucun couple $x,y\in A$ pour lequel on ait à la fois $Rxy$ et~$Ryx$.
\end{defn}

Si $A \neq \emptyset$, aucune relation irréflexive sur~$A$ n’est réflexive,
et toute relation asymétrique sur~$A$ est aussi antisymétrique.
Il existe cependant, dès que $A$ contient au moins deux éléments,
des relations $R \subseteq A^2$ qui ne sont ni réflexives ni
irréflexives ; et, sur tout ensemble non vide, il existe des relations
antisymétriques qui ne sont pas asymétriques.\footnote{Précision éditoriale :
la première affirmation d’existence de l’original demande au moins deux
éléments. Sur un singleton, les seules relations sont la relation vide,
irréflexive, et la relation d’identité, réflexive.}

\end{document}
